\begin{abstract}
This paper presents MARS, a terminal editor, multiplexer, and LLM agent that ship as a
single binary. Because the agent is embedded where the user works, it can accumulate
two sources of context that generic assistants see only in fragments and never keep:
a \emph{reflex} memory of the user's own accepted commands, and a \emph{reflective}
memory of the tool's own documentation and configuration registry. MARS retrieves from both with plain lexical search (BM25 over
local files; no embeddings, no data leaving the machine) and injects the results into
the model's prompt. In the constrained domain of the terminal this simple
mechanism is unusually effective. Under a leakage-controlled protocol in which the
memory store holds only \emph{paraphrases} of the user's earlier requests, retrieval
raises the accuracy of a small, inexpensive model (\code{claude-haiku-4-5}) on
project-specific commands from 38\% (9/24) to 96\% (23/24), so a win requires the
retriever to bridge a paraphrase gap rather than look up a planted answer. Pointed at
the tool's own documentation, the same mechanism raises accuracy on questions about
the tool itself from 43\% to 76\% overall, with the decisive gain on plain-English
reconfiguration, from 17\% to 93\%. We describe the registry-grounded directive protocol, the corrective-memory loop,
the architecture that keeps the agent present across machines while the API key stays
home, and a self-instrumenting evaluation built from the agent's own debug log. MARS is open
source (MIT), installs with \code{cargo install mars-terminal}, and ships with a
screencast.
\end{abstract}
